% GENERATED FILE. Edit canonical lessons, then run npm run generate:books.
% canonical-source-hash: fnv1a64:a9a4509fead385c3
% canonical-lessons: FR-C457-clairiere, FR-C457-buisson, FR-C457-fougere, FR-C457-mousse, FR-C457-ecorce

\chapter{Au fond des bois}
\label{ch:fr-au-fond-des-bois}
\hlchaptermodality{457}

\begin{chapteropening}
\textbf{By the end of this chapter:} \emph{I can describe a wood: the clearing, the bush, the fern, the moss and the bark.}

Complete the last lesson of chapter 457: I can describe a wood: the clearing, the bush, the fern, the moss and the bark.
\end{chapteropening}

\section[la clairière]{la clairière — the clearing (in a wood)}

\label{lesson:FR-C457-clairiere}



\begin{quote}
Before the new one: say the French for to splash about, then the French for to stream.
\end{quote}

\subsection*{You'll want to know: la clairière}
\textbf{la clairière} — \textquotedblleft{}the clearing (in a wood)\textquotedblright{}.

\begin{quote}
Au milieu de la forêt, nous avons pique-niqué dans une clairière.
\end{quote}

\begin{quote}
\emph{In the middle of the forest we had a picnic in a clearing.}
\end{quote}

Built on \emph{clair} 'light'.

\begin{cousinweb}
\emph{clairière} is built on \emph{clair} 'clear, light', which is inherited from Latin \emph{clārus} (Lewis \& Short): the light place among the trees. EtymDB does not list it.
\end{cousinweb}

\begin{grammarlens}[title={What you've built}]
One more word for telling, explaining and describing at B1.
\end{grammarlens}

\subsection*{Your turn}
\begin{itemize}
\raggedright
  \item \emph{Say it:} \emph{la clairière}
  \item \emph{Say it:} \emph{la clairière}, once more
  \item \emph{Say it:} a sentence of your own with \emph{la clairière}
  \item \emph{Recall:} say the French for to splash about, then the French for to stream, then say \emph{la clairière} again
\end{itemize}

\begin{tcolorbox}[breakable,colback=teal!4,colframe=teal!35!black,title={Before you move on}]
What does la clairière mean? (\textquotedblleft{}The clearing (in a wood)\textquotedblright{}.) Say it once more.
\end{tcolorbox}
\section[le buisson]{le buisson — the bush}

\label{lesson:FR-C457-buisson}



\begin{quote}
Before the new one: say the French for to stream, then the French for the clearing.
\end{quote}

\subsection*{You'll want to know: le buisson}
\textbf{le buisson} — \textquotedblleft{}the bush\textquotedblright{}.

\begin{quote}
Un lapin se cache dans le buisson.
\end{quote}

\begin{quote}
\emph{A rabbit is hiding in the bush.}
\end{quote}

\emph{Tourner autour du pot} means to beat about the bush.

\begin{cousinweb}
\emph{buisson} is built on \emph{bois} 'wood', which is usually traced to a Germanic word for 'bush' (EtymDB lists Germanic cousins).
\end{cousinweb}

\begin{grammarlens}[title={What you've built}]
One more word for telling, explaining and describing at B1.
\end{grammarlens}

\subsection*{Your turn}
\begin{itemize}
\raggedright
  \item \emph{Say it:} \emph{le buisson}
  \item \emph{Say it:} \emph{le buisson}, once more
  \item \emph{Say it:} a sentence of your own with \emph{le buisson}
  \item \emph{Recall:} say the French for to stream, then the French for the clearing, then say \emph{le buisson} again
\end{itemize}

\begin{tcolorbox}[breakable,colback=teal!4,colframe=teal!35!black,title={Before you move on}]
What does le buisson mean? (\textquotedblleft{}The bush\textquotedblright{}.) Say it once more.
\end{tcolorbox}
\section[la fougère]{la fougère — the fern, the bracken}

\label{lesson:FR-C457-fougere}



\begin{quote}
Before the new one: say the French for the clearing, then the French for the bush.
\end{quote}

\subsection*{You'll want to know: la fougère}
\textbf{la fougère} — \textquotedblleft{}the fern, the bracken\textquotedblright{}.

\begin{quote}
Les fougères poussent sous les arbres.
\end{quote}

\begin{quote}
\emph{The ferns grow under the trees.}
\end{quote}

Often in the plural.

\begin{cousinweb}
\emph{fougère} comes from Old French \emph{fulgiere}, from Vulgar Latin \emph{filicāria} 'place full of ferns', from Latin \emph{filix} 'fern' (EtymDB).
\end{cousinweb}

\begin{grammarlens}[title={What you've built}]
One more word for telling, explaining and describing at B1.
\end{grammarlens}

\subsection*{Your turn}
\begin{itemize}
\raggedright
  \item \emph{Say it:} \emph{la fougère}
  \item \emph{Say it:} \emph{la fougère}, once more
  \item \emph{Say it:} a sentence of your own with \emph{la fougère}
  \item \emph{Recall:} say the French for the clearing, then the French for the bush, then say \emph{la fougère} again
\end{itemize}

\begin{tcolorbox}[breakable,colback=teal!4,colframe=teal!35!black,title={Before you move on}]
What does la fougère mean? (\textquotedblleft{}The fern, the bracken\textquotedblright{}.) Say it once more.
\end{tcolorbox}
\section[la mousse]{la mousse — the moss; the foam}

\label{lesson:FR-C457-mousse}



\begin{quote}
Before the new one: say the French for the bush, then the French for the fern.
\end{quote}

\subsection*{You'll want to know: la mousse}
\textbf{la mousse} — \textquotedblleft{}the moss; the foam\textquotedblright{}.

\begin{quote}
Les pierres du chemin sont couvertes de mousse.
\end{quote}

\begin{quote}
\emph{The stones on the path are covered in moss.}
\end{quote}

\emph{La mousse} is also foam, as on beer; that is probably a different word.

\begin{cousinweb}
\emph{mousse} 'moss' comes from Old French \emph{mosse}, from Frankish \emph{mosa} 'moss' (EtymDB).
\end{cousinweb}

\begin{grammarlens}[title={What you've built}]
One more word for telling, explaining and describing at B1.
\end{grammarlens}

\subsection*{Your turn}
\begin{itemize}
\raggedright
  \item \emph{Say it:} \emph{la mousse}
  \item \emph{Say it:} \emph{la mousse}, once more
  \item \emph{Say it:} a sentence of your own with \emph{la mousse}
  \item \emph{Recall:} say the French for the bush, then the French for the fern, then say \emph{la mousse} again
\end{itemize}

\begin{tcolorbox}[breakable,colback=teal!4,colframe=teal!35!black,title={Before you move on}]
What does la mousse mean? (\textquotedblleft{}The moss; the foam\textquotedblright{}.) Say it once more.
\end{tcolorbox}
\section[l'écorce]{l'écorce — the bark (of a tree); the peel}

\label{lesson:FR-C457-ecorce}



\begin{quote}
Before the new one: say the French for the fern, then the French for the moss.
\end{quote}

\subsection*{You'll want to know: l'écorce}
\textbf{l'écorce} — \textquotedblleft{}the bark (of a tree); the peel\textquotedblright{}.

\begin{quote}
L'écorce de ce vieil arbre est très épaisse.
\end{quote}

\begin{quote}
\emph{The bark of this old tree is very thick.}
\end{quote}

A feminine noun. \emph{L'écorce d'orange} — orange peel.

\begin{cousinweb}
\emph{écorce} is inherited from Latin \emph{scortea} 'leather garment', from \emph{scortum} 'hide, skin' (EtymDB): the tree's skin.
\end{cousinweb}

\begin{grammarlens}[title={What you've built}]
One more word for telling, explaining and describing at B1.
\end{grammarlens}

\subsection*{Your turn}
\begin{itemize}
\raggedright
  \item \emph{Say it:} \emph{l'écorce}
  \item \emph{Say it:} \emph{l'écorce}, once more
  \item \emph{Say it:} a sentence of your own with \emph{l'écorce}
  \item \emph{Recall:} say the French for the fern, then the French for the moss, then say \emph{l'écorce} again
\end{itemize}

\begin{tcolorbox}[breakable,colback=teal!4,colframe=teal!35!black,title={Before you move on}]
What does l'écorce mean? (\textquotedblleft{}The bark (of a tree); the peel\textquotedblright{}.) Say it once more.
\end{tcolorbox}
